\documentclass[twocolumn]{aastex631}
\usepackage{amsmath}
\usepackage{booktabs}
\shorttitle{SDSS Density Proxy for Environmental Quenching}
\shortauthors{NebulaMind Research Autopilot}
\begin{document}

\title{SDSS Density Proxy for Environmental Quenching}
\author{NebulaMind Research Autopilot}
\affiliation{NebulaMind Astrophysics Collaboration, San Francisco, CA 94107, USA}
\correspondingauthor{NebulaMind Research Autopilot}
\email{autopilot@nebulamind.ai}

\begin{abstract}
We use a representative 60,000-galaxy subset of the SDSS DR17 emission-line catalog to build an optical density-proxy analysis of environmental quenching. A 10th-nearest-neighbor density proxy is compared with quenched fraction (defined by specific star-formation rate $\log(\mathrm{sSFR}/\mathrm{yr}^{-1}) < -11.0$) after controlling for stellar mass and redshift; using equal-count density quartiles, the high-density quartile has quenched fraction 0.230 $\pm$ 0.003 versus 0.181 $\pm$ 0.003 in the low-density quartile. The bootstrap high-minus-low interval is [0.041, 0.059], which excludes zero. This analysis is intentionally limited to the optical denominator and leaves the missing group and halo information for future study.
\end{abstract}

\keywords{galaxies: evolution --- galaxies: active --- galaxies: star formation --- surveys --- methods: data analysis}
\software{Astropy, SciPy, NumPy, Matplotlib, pandas}

\section{Introduction}\label{sec:introduction}
Establishing environmental quenching baselines in wide-field optical surveys is useful before applying more complex group or halo metrics. Here, the SDSS DR17 emission-line parent sample serves as the optical baseline for the environmental-quenching comparison rather than a direct physical measurement. We evaluate a local 10th-nearest-neighbor density proxy using SDSS DR17 emission-line galaxies and restrict the scope to directly measured optical properties. Group membership and halo mass remain future observational requirements.


\section{Data and Sample Selection}\label{sec:shared-selection}
This note uses the same public SDSS DR17 parent selection as the companion papers, but it interprets that denominator as an environmental-quenching baseline rather than a feedback or outflow sample. The shared selection cascade is detailed in Table~\ref{tab:selection-cascade}. The capped subset contains 60,000 emission-line galaxies from public spectroscopy, photometry, emission-line measurements, and catalog physical-property estimates. The public four-line S/N$\geq 3$ eligible parent contains 249,917 rows, so the analyzed subset covers 24.0\% of that parent. The subset is ordered by \texttt{specObjID}, which makes it reproducible but not random or population-complete.

\begin{deluxetable*}{lrrr}
\tabletypesize{\scriptsize}
\tablecaption{Shared SDSS DR17 selection cascade used before paper-specific quantities.\label{tab:selection-cascade}}
\tablehead{\colhead{Selection stage} & \colhead{Public DR17 rows} & \colhead{Cached rows} & \colhead{Retention vs. spectro-z parent (fraction)}}
\startdata
SpecObj GALAXY, $0.02 < z < 0.12$ & 501,060 & -- & 1.000 \\
plus galSpecInfo/PhotoObj/galSpecExtra and mass/sSFR bounds & 416,554 & -- & 0.831 \\
plus galSpecLine join & 416,554 & -- & 0.831 \\
four BPT lines positive with positive errors & 373,445 & 60,000 & 0.745 \\
four BPT lines S/N$\geq 3$ & 249,917 & 60,000 & 0.499 \\
four BPT lines S/N$\geq 5$ & 176,523 & 42,446 & 0.352 \\
four BPT lines S/N$\geq 10$ & 91,768 & 22,311 & 0.183 \\
\enddata
\tablecomments{Counts are read-only public SDSS DR17 count queries plus the cached local CSV. Cached rows are shown only where the cache applies.}
\end{deluxetable*}

The four-line requirement is strongly selection dependent. In the public counts, S/N$\geq 3$ keeps 33.6\% of the $-12<\log(\mathrm{sSFR}/\mathrm{yr}^{-1})<-11$ parent bin but 94.9\% of the $-10<\log(\mathrm{sSFR}/\mathrm{yr}^{-1})<-9.5$ bin. Therefore every incidence, quenching, density, gas-denominator, or target-vector statement in this analysis is conditional on the four-line emission-line selection.

Local subset versus public catalog marginal checks found no redshift, stellar-mass, or sSFR bin with a subset-minus-public fraction difference above 5 percentage points. The largest absolute differences were 2.03 percentage points in redshift, -1.63 percentage points in stellar mass, and -0.58 percentage points in sSFR. This is a representativeness diagnostic only; it does not make the capped cache random or complete.


\section{Measurements}\label{sec:measurements}
The row-level measurements include redshift, stellar mass, catalog specific star-formation rate, model colors, and four optical line fluxes/errors. BPT labels and all derived denominators are recomputed locally from the cached table \citep{baldwin1981,kewley2001,kauffmann2003bpt,kewley2006}. Catalog SFR/sSFR values are treated as low-redshift SDSS physical-property estimates rather than direct resolved gas or feedback measurements \citep{sdssdr17,brinchmann2004,york2000}.
Unless otherwise noted, quoted fraction uncertainties are binomial counting uncertainties from the stated sample sizes, and bracketed intervals are bootstrap confidence intervals.


\section{SDSS density-proxy result for environmental quenching}\label{sec:topic-result}
We examine whether a nearest-neighbor density proxy adds quenched-fraction information beyond stellar mass in the SDSS emission-line sample. The result is an optical baseline rather than a full physical-feedback test.

Within the 60,000-galaxy denominator, the high-density quartile has a quenched fraction of $0.230 \pm 0.003$ ($n=3{,}456$ of $15{,}000$). The low-density quartile is lower at $0.181 \pm 0.003$ ($n=2{,}710$ of $15{,}000$). The bootstrap high-minus-low difference interval is $[0.041, 0.059]$, safely excluding zero. After controlling for logarithmic stellar mass and redshift, a linear probability model yields a high-density coefficient of $0.032 \pm 0.004$. That result indicates the density proxy remains correlated with quenching independently of those host-galaxy properties. Figure~\ref{fig:topic} summarizes the proxy split used in this baseline.


\begin{figure}
\centering
\includegraphics[width=\columnwidth]{../figures/fig-topic.pdf}
\caption{SDSS DR17 optical density-proxy diagnostic for environmental quenching. The figure shows quenched fraction versus the local density proxy, split into equal-count density quartiles. The high-density quartile reaches a quenched fraction of 0.230 $\pm$ 0.003, establishing the baseline for future group-catalog analyses.}
\label{fig:topic}
\end{figure}

\section{Interpretation and missing observables}\label{sec:missing}
This SDSS-only baseline does not include group catalogues, robust central/satellite labels, halo masses, morphology, or multi-redshift selection functions. Those data are required before the density proxy can be interpreted as a physical environmental-quenching measurement.

Because the density proxy is projected on the sky, fiber-collision and redshift-space incompleteness can dilute nearest-neighbor estimates; future group-catalog work should correct those effects explicitly.

Mass and environment are separable axes in low-redshift galaxy evolution, and prior studies have developed that framework \citep{peng2010,baldry2006,wetzel2013,goubert2024}. A complete environmental-quenching analysis therefore requires group and halo information beyond this nearest-neighbor proxy, together with central/satellite labels.


\section{Data Availability}\label{sec:data-avail}
The SDSS DR17 spectroscopy, photometry, emission-line measurements, and catalog quantities used here are publicly available through the SDSS data releases and associated catalog papers cited below. A local subset and manifest are retained in the project repository for reproducibility and are available from the corresponding author upon reasonable request.

\section{Conclusion}\label{sec:conclusion}
The SDSS-only proxy shows a high-density quenched fraction of 0.230 $\pm$ 0.003 versus 0.181 $\pm$ 0.003 in the low-density quartile, with a mass- and redshift-adjusted high-density coefficient of $0.032 \pm 0.004$. These values define an optical environmental baseline, but a full quenching interpretation still requires group catalogs, halo masses, and central/satellite labels.

\begin{acknowledgments}
We thank the SDSS collaboration. This manuscript uses public SDSS DR17 data only.
\end{acknowledgments}
\begin{thebibliography}{99}
\bibitem[Abdurro'uf et al.(2022)]{sdssdr17} Abdurro'uf, Accetta, K., Aerts, C., et al. 2022, ApJS, 259, 35
\bibitem[Baldwin et al.(1981)]{baldwin1981} Baldwin, J.~A., Phillips, M.~M., \& Terlevich, R. 1981, PASP, 93, 5
\bibitem[Brinchmann et al.(2004)]{brinchmann2004} Brinchmann, J., Charlot, S., White, S.~D.~M., et al. 2004, MNRAS, 351, 1151
\bibitem[Kauffmann et al.(2003a)]{kauffmann2003bpt} Kauffmann, G., Heckman, T.~M., Tremonti, C., et al. 2003a, MNRAS, 346, 1055
\bibitem[Kewley et al.(2001)]{kewley2001} Kewley, L.~J., Dopita, M.~A., Sutherland, R.~S., Heisler, C.~A., \& Trevena, J. 2001, ApJ, 556, 121
\bibitem[Kewley et al.(2006)]{kewley2006} Kewley, L.~J., Groves, B., Kauffmann, G., \& Heckman, T. 2006, MNRAS, 372, 961
\bibitem[York et al.(2000)]{york2000} York, D.~G., Adelman, J., Anderson, J.~E., Jr., et al. 2000, AJ, 120, 1579
\bibitem[Baldry et al.(2006)]{baldry2006} Baldry, I.~K., Balogh, M.~L., Bower, R.~G., et al. 2006, MNRAS, 373, 469
\bibitem[Goubert et al.(2024)]{goubert2024} Goubert, P.~H., Bluck, A.~F.~L., Piotrowska, J.~M., \& Maiolino, R. 2024, MNRAS, 528, 3822
\bibitem[Peng et al.(2010)]{peng2010} Peng, Y.-j., Lilly, S.~J., Kovac, K., et al. 2010, ApJ, 721, 193
\bibitem[Wetzel et al.(2013)]{wetzel2013} Wetzel, A.~R., Tinker, J.~L., Conroy, C., \& van den Bosch, F.~C. 2013, MNRAS, 432, 336
\end{thebibliography}

\end{document}
